\documentclass[11pt,a4paper,sans]{moderncv}
\moderncvstyle{banking}
\moderncvcolor{black} 
\usepackage[scale=0.88]{geometry}
\usepackage{enumitem}
\pagenumbering{arabic}
\usepackage{hyperref}

\hypersetup{colorlinks=true, urlcolor=[rgb]{0.15,0.15,0.45}}
\setlist[itemize]{noitemsep,topsep=2pt,leftmargin=1.5em}
\renewcommand*{\namefont}{\fontsize{18}{20}\selectfont}
\renewcommand*{\titlefont}{\fontsize{14}{16}\selectfont}

% ----- personal data -----
\name{Sushant}{Gautam}
\title{Agentic AI and Multimodal Understanding}
\phone[mobile]{+47 986 72 556}
\email{sushant@simula.no}
\homepage{sushant.info.np}
\social[linkedin]{eSushant}
\social[github]{SushantGautam}
\social[researchgate]{Sushant-Gautam}
\social[googlescholar]{t3Iie8cAAAAJ}

\begin{document}
\makecvtitle
\vspace*{-2.0\baselineskip}

%========================
\section{Profile}
\cvitem{}{
AI engineer with 5+ years of experience building and deploying machine learning systems across startup and research environments. Currently pursuing a PhD in agentic AI and multimodal understanding with focus on robust and trustworthy systems. Experience includes developing production ML pipelines, APIs, and real-time AI applications, along with research on vision–language models and intelligent agents. Strong interest in LLM-based system evaluation, and scalable AI infrastructure.
}

%========================
\section{Work Experience}

\cventry{May 2023 -- Present}
{Ph.D. Student (Full-time)}
{\href{https://www.simulamet.no/people/sushant}{Simula Metropolitan Center for Digital Engineering (SimulaMet)}}
{Oslo, Norway}
{Affiliated with \href{https://www.oslomet.no/en}{Oslo Metropolitan University}}
{
\begin{itemize}
\item Conduct research on multimodal AI systems with focus on robustness, uncertainty, and real-world applicability.
\item Work on vision–language models and intelligent systems involving multiple modalities and reasoning.
\item Explore system-level behavior of AI models, including reliability and limitations in practical settings.
\item Supervise MSc students working on applied AI and machine learning projects.
\item Collaborate across research and engineering teams to translate research ideas into practical systems.
\end{itemize}
}

\cventry{Oct 2024 -- Jul 2025}
{Ph.D. Enrichment Student for the Turing Enrichment Scheme 2024/25}
{\href{https://www.turing.ac.uk/people/doctoral-students/sushant-gautam}{The Alan Turing Institute}}
{London, United Kingdom}
{}
{
\begin{itemize}
\item Contributed to development of multimodal datasets and benchmarks for medical reasoning.
\item Participated in interdisciplinary collaborations across research labs and applied AI teams.
\item Engaged in industry-oriented sessions (e.g., Turing–Roche) focusing on real-world AI applications.
\end{itemize}
}

\cventry{Oct 2022 -- May 2023}
{Adjunct Lecturer (Contract-based)}
{\href{https://tcioe.edu.np/}{Thapathali Campus}, Tribhuvan University}
{Kathmandu, Nepal}
{}
{
\begin{itemize}
\item Delivered lectures and labs on applied deep learning.
\item Supervised student projects focused on real-world ML applications.
\end{itemize}
}

\cventry{Apr 2019 -- May 2023}
{Joined as AI Engineer and later promoted to Engineering Manager}
{\href{https://www.nsdevil.com/?lang=en}{North Star Developer's Village}}
{Lalitpur, Nepal}
{}
{
\begin{itemize}
\item Led development and deployment of AI-powered systems in production environments.
\item Built ML-based face verification and exam integrity systems with real-time inference.
\item Designed and deployed backend APIs for NLP systems including essay scoring and feedback generation.
\item Developed scalable ML pipelines for training, evaluation, and deployment.
\item Managed engineering efforts balancing model performance, system reliability, and product requirements.
\end{itemize}
}

\cventry{Oct 2018 -- Mar 2019}
{AI Intern}
{\href{https://www.lftechnology.com/}{Leapfrog Technology}}
{Kathmandu, Nepal}
{}
{
\begin{itemize}
\item Developed time-series prediction models for weather and pollution forecasting.
\item Deployed models with web-based interfaces for interactive usage.
\end{itemize}
}

%========================
\section{Technical Capabilities}

\cvitem{Programming}{Python (primary), Java/Kotlin, SQL, Bash, C/C++.}

\cvitem{AI/ML}{PyTorch, TensorFlow/Keras, HuggingFace, Scikit-learn, OpenCV, PEFT/LoRA.}

\cvitem{LLM \& Systems}{LangChain, vLLM, RAG pipelines, multimodal systems, model evaluation workflows.}

\cvitem{Backend \& Infra}{FastAPI, Django, REST APIs, Docker, Git, CI/CD workflows.}

\cvitem{Data \& Tools}{NumPy, Pandas, Matplotlib, Jupyter, data pipelines.}

\cvitem{Languages}{English (Fluent), Nepali (Native), Hindi (Advanced), Norwegian (Basic).}

%========================
\section{Selected Projects \& Research Highlights}

\cvitem{\textbf{SimpleAudit (AI Safety \& Evaluation Framework)}}{
\href{https://www.simula.no/research/projects/simpleaudit}{Project Link} \\
Contributed to development of a framework for evaluating large language models and AI systems in real-world scenarios. Focus on identifying failure modes, robustness issues, and safety risks beyond standard benchmarks. Supports systematic testing of AI behavior under realistic conditions.
}

\cvitem{\textbf{MoltBook Observatory (Multi-Agent AI System)}}{
\href{https://www.forskning.no/data-informasjonsteknologi-kunstig-intelligens/pa-dette-nettstedet-snakker-ki-fritt-med-hverandre-lagde-egen-religion/2623586}{Article Link} \\
Developed and analyzed an experimental multi-agent platform where AI agents interact autonomously in a shared environment. Studied emergent behaviors, content dynamics, and system-level properties such as novelty, repetition, and interaction patterns. Demonstrates understanding of agent orchestration and large-scale AI system behavior.
}

\cvitem{\textbf{Kvasir-VQA-x1 (Multimodal AI Benchmark)}}{
\href{https://github.com/Simula/Kvasir-VQA-x1}{Project Link} \\
Designed a large-scale multimodal benchmark with 150k+ question–answer pairs for medical visual question answering. Incorporated LLM-generated reasoning tasks with varying complexity and introduced robustness evaluation through realistic perturbations.
}

\cvitem{\textbf{Kvasir-VQA (Vision--Language Dataset)}}{
\href{https://datasets.simula.no/kvasir-vqa}{Dataset Link} \\
Built a multimodal dataset of 6,500+ annotated medical images supporting VQA, captioning, and classification tasks for real-world diagnostic scenarios.
}

\cvitem{\textbf{AI-Based Soccer Game Summarization (Multimodal System)}}{
\href{https://www.researchgate.net/publication/363857936_AI-based_Soccer_Game_Summarization_From_Video_Highlights_to_Dynamic_Text_Summaries}{Project Link} \\
Developed an end-to-end system integrating audio, video, and metadata to generate structured natural language summaries. Demonstrates multimodal pipeline design and applied AI system development.
}

\cvitem{\textbf{Pose-Based Dystonia Score Prediction (Applied Computer Vision)}}{
\href{https://github.com/SushantGautam/DystViz/}{Project Link} \\
Built a deep learning system using pose estimation for clinical assessment of dystonia severity. Applied computer vision in a healthcare context with focus on robustness and interpretability.
}

%========================
\section{Education}

\cventry{May 2023 -- Jul 2026 (Expected)}
{Doctor of Philosophy (Ph.D.), Computer Engineering}
{\href{https://www.oslomet.no/en}{Oslo Metropolitan University (OsloMet)}}
{Oslo, Norway}
{}
{
\textbf{Thesis:} Bridging Multimedia Modalities: Enhanced Multimodal AI Understanding and Intelligent Agents.\\
\textbf{Focus:} Multimodal AI, vision--language models, intelligent agents, uncertainty and robustness in AI systems.\\
\textbf{Coursework:} Topics in AI and Machine Learning; Scientific Research Methods and Data Analysis; Engineering Science and Ethics.\\
\textbf{Additional Coursework:} Neural Methods in Natural Language Processing (\href{https://www.uio.no/}{University of Oslo}).
}

\cventry{Jan 2024 -- Jul 2025}
{Ph.D. Enrichment Scheme Placement}
{\href{https://www.turing.ac.uk/}{The Alan Turing Institute}}
{London, United Kingdom}
{}
{
Selected through NORA for international research training; contributed to multimodal AI benchmarks and interdisciplinary collaboration with academic and industry partners.
}

\cventry{Oct 2020 -- Sep 2022}
{Master of Engineering (MEng), Informatics \& Intelligent Systems}
{\href{https://tcioe.edu.np/}{Tribhuvan University, Thapathali Campus}}
{Nepal}
{}
{
\textit{Full scholarship for academic merit; top of class.}\\
\textbf{Relevant Coursework:} Deep Learning, Computer Vision, Natural Language Processing, Applied Machine Learning, Optimization, Information Theory.\\
\textbf{Thesis:} AI-Based Soccer Game Summarization: From Video Highlights to Dynamic Text Summaries.
}

\cventry{Oct 2015 -- Oct 2019}
{Bachelor of Engineering (BE), Computer Engineering}
{\href{https://pcampus.edu.np/}{Tribhuvan University, Pulchowk Campus}}
{Nepal}
{}
{
\textit{Full government scholarship for national-level entrance ranking.}\\
\textbf{Relevant Coursework:} Artificial Intelligence, Data Mining, Probability \& Statistics, Big Data.\\
\textbf{Projects:} Weather and pollution prediction using RNN and ARIMA; facial landmark detection using neural networks.
}


%========================
\section{References}
Available upon request.

\vfill
\begin{flushright}
{\footnotesize Compiled: \today}
\end{flushright}

\end{document}